% Fichier généré automatiquement par tools/generate_corrections.py.
% Source : DYN/DYN-01/61_Hemostase/61_Hemostase.tex
% Statut : partiel (6/7 question(s) renseignée(s), 0 reprise(s)).
\begin{corrigebox}[Corrigé partiel extrait de la banque]
\CorrigeQuestion{1}
La distance  $x_M^{\text{max}}$ correspond à l'aire sous le courbe de la loi de commande de vitesse.
On a alors 
 $x_M^{\text{max}} = \left(T-T_a\right)V_M^x$ 
 $ \Longleftrightarrow V_M^x=\dfrac{x_M^{\text{max}}}{T-T_a}$.
\CorrigeQuestion{2}
\begin{itemize}
\item Expression de l'énerige cinétique : $\ec{E}{0}=\dfrac{1}{2}J_e (\omega_m^x)^2$.
\item Puissance intérieure : $\mathcal{P}_{\text{int}}(E)=0$.
\item Puissance extérieure : $\mathcal{P}_{\text{ext}}(E)=C_m\omega_m^x$.
\item Application du TEC : $J_e \omega_m^x \dot{\omega}_m^x = C_m\omega_m^x$ soit $J_e  \dot{\omega}_m^x = C_m$.
\end{itemize}
\vspace{.25cm}

On a alors sur chacune des phases : 

\begin{itemize}
\item Phase 1 : $C_m=J_e  \dot{\omega}_m^x$ avec $\dot{\omega}_m^x = \dot{V_M}^x /  \lambda= \dfrac{V_M^x}{ \lambda T_a}$ et $C_m = J_e \dfrac{V_M^x}{\lambda T_a}$.
\item Phase 2 : $C_m=0$.
\item Phase 3 : $C_m = - J_e  \dfrac{V_M^x}{\lambda T_a}$.
\end{itemize}
\CorrigeQuestion{3}
\CorrigeACompleter{Aucun élément de réponse n'est présent dans la source d'origine pour cette question.}
\CorrigeQuestion{4}
$P_{\text{max}} = J_e \dfrac{V_M^x}{\lambda T_a}  {\omega}_m^x =J_e \dfrac{(V_M^x)^2}{\lambda^2 T_a}$.
\CorrigeQuestion{5}
On a alors 
$P_{\text{max}} =J_e \dfrac{\left({x_M^{\text{max}}}\right)^2}{\lambda^2 \left({T-T_a}\right)^2 T_a}$.
\CorrigeQuestion{6}
On résout $\dfrac{\dd P_{\text{max}}}{\dd T_a} = 0$ et on cherche la valeur de $T_a$ pour laquelle 
$P_{\text{max}}$ est minimale.
\CorrigeQuestion{7}
On a  $V_M^x=\dfrac{x_M^{\text{max}}}{T-T_a}$. D'autre part, 
$V_M^x = \omega_M^x  k R_p $ soit $\omega_M^x  = \dfrac{V_M^x}{ k R_p} =  \dfrac{{x_M^{\text{max}}}}{ k R_p\left({T-T_a} \right)} $.

AN : $\omega_M^x  = \dfrac{550\times 10^{-3}}{ 0,1 \times 20 \times 10^{-3} \left(1-1/3 \right)} $
$ =  \dfrac{550\times 3}{ 4}= \SI{412,5}{rad.s^{-1}} $ soit \SI{3941}{tr.min^{-1}}.

Cette valeur est bien compatible avec la vitesse du moteur.
\end{corrigebox}
